\section{Introduction}
The path-ordered phase factor $\text{P}\exp\left(ig\int A_{\mu} dx^{\mu}\right)$, referred to as the Wilson line\footnote{Throughout this work, our focus will be exclusively on off-lightcone Wilson lines and Wilson loops, which we will refer to simply as Wilson lines unless ambiguity arises.}, is  a fundamental object  in gauge theories,  playing a crucial role in connecting local fields to construct gauge-invariant operators. The study of the renormalization of  Wilson lines and Wilson loops was initiated by Polyakov \cite{Polyakov:1980ca}  and Gervais, Neveu \cite{Gervais:1979fv}. 
The all-order proof of the multiplicative renormalizability of Wilson lines was made using diagrammatic methods in refs.~\cite{Dotsenko:1979wb,Craigie:1980qs} and functional approach in ref.~\cite{Dorn:1986dt} based on the one-dimensional auxiliary-field formalism \cite{Gervais:1979fv,Arefeva:1980zd, Arefeva:1980uy,Samuel:1978iy,Brandt:1978wc}; see ref.~\cite{Dorn:1986dt} for a review. 

Off-lightcone Wilson-line operators are constructed using local operators connected by time-like or space-like Wilson lines, which insure the gauge invariance of the non-local operators.  Wilson-line operators have broad applications in various contexts. For instance, space-like Wilson-line operators play a crucial role in determining quasi-distribution function (quasi-PDFs) \cite{Ji:2013dva} and pseudo-distribution functions (pseudo-PDFs) \cite{Radyushkin:2017cyf}, while time-like Wilson-line operators are essential for understanding quarkonium decay and production within the potential non-relativistic QCD (pNRQCD) effective theory \cite{Pineda:1997bj,Brambilla:1999xf,Brambilla:2004jw,Brambilla:2001xy,Brambilla:2002nu,Brambilla:2020ojz,Brambilla:2021abf,Brambilla:2022rjd,Brambilla:2022ayc}.

Recently, there has been increasing interests in the study of the quasi-PDFs and pseudo-PDFs defined as the hadronic matrix elements of the Wilson-line operators of the type
\begin{eqnarray}
\label{eq:defqoperator}
\mathcal{O}_{\Gamma}(zv) &=& \bar{\psi}(zv)\Gamma W(zv,0)\psi(0),\\
\label{eq:defgoperator}
\mathcal{O}^{\mu\nu\alpha\beta} (zv)&=& g^2 F^{\mu\nu}(zv)W(zv,0)F^{\alpha\beta}(0),
\end{eqnarray}
where $\psi$ is a quark field, $F^{\mu\nu}$ is the gluon field strength tensor, $\Gamma$ represents a Dirac structure, $v$ is a four-vector satisfying $v^2=-1$ (or $v^2=1$) for space-like (or time-like) Wilson-line operators in Minkowski spacetime \footnote{In this context and throughout subsequent discussions, we stay in Minkowski spacetime. However, a transition to Euclidean spacetime  is warranted when introducing gradient flow, as it is defined in Euclidean spacetime, which indicates $v^2=1$ both for space-like and time-like Wilson-line operators.}, $z$ is a real number and $W(zv, 0)$ is a straight-line path-ordered Wilson line connecting the two fields,
\begin{eqnarray}
W(zv, 0) = \text{P}\exp\left(ig\int_0^z ds\, v\cdot A(sv) \right),
\end{eqnarray}
in which the fundamental representation and the adjoint representation of the gauge group are assumed for the quark case, and the gluon case, respectively. Without loss of generality, we choose $z>0$ throughout this work. Notably,  quasi-PDFs and pseudo-PDFs can be factorized in relation to the standard light-cone parton distribution functions (PDFs) based on the short-distance operator product expansion (OPE) and the large momentum effective theory (LaMET) \cite{Braun:2007wv, Ji:2014gla,Izubuchi:2018srq}.  Employing hadronic matrix elements of  off-lightcone Wilson-line operators as opposed to lightlike ones, offers the advantage of being considerably more suitable to lattice computations. However, this convenience comes with the trade-off of the highly non-trivial renormalization of the off-lightcone Wilson-line operators. For a comprehensive overview, we refer to ref.~\cite{Ji:2020ect}.
Furthermore, the vacuum expectation value of gluonic Wilson-line operators, such as the one defined in eq.~(\ref{eq:defgoperator}), are related to the studies of quarkonium decay and production in the framework of pNRQCD, and the QCD vacuum structure 
(see ref.~\cite{DiGiacomo:2000irz} for a review).


The gradient flow was originally introduced as  a method to regularize and suppress the ultraviolet (UV) fluctuations \cite{Narayanan:2006rf,Luscher:2009eq} in lattice calculations. It has proven useful in lattice QCD scale settings and in lattice calculation of local matrix elements and correlation functions \cite{Luscher:2010iy,Luscher:2011bx,Luscher:2013cpa,Luscher:2013vga,Borsanyi:2012zs,Sommer:2014mea,Suzuki:2013gza,Makino:2014taa,Harlander:2018zpi,FlavourLatticeAveragingGroupFLAG:2021npn,Leino:2021vop,Mayer-Steudte:2022uih,Leino:2022kgj}. Alongside its noise-reducing effect, gradient flow exhibits the notable property that composite operators at positive flow time remain finite even in the continue limit,  provided they are expressed using renormalized physical parameters and fields \cite{Luscher:2011bx}. This property renders gradient flow suitable as a matching scheme between lattice and perturbative calculations. It is worth noting that gradient flow does not modify the infrared behavior of operators in the small flow-time limit, allowing the use of dimensional regularization to regulate infrared (IR) divergences when matching from the gradient-flow scheme to $\overline{\rm MS}$-like schemes.

The initial proposal to explore quasi-PDFs with the incorporation of gradient flow was presented in ref.~\cite{Monahan:2016bvm}. Subsequently, a one-loop matching calculation was performed, considering full flow time dependencies, for the hadronic matrix element defined using the operator in eq.~(\ref{eq:defqoperator}) with zero external momenta, which is related to the quark quasi-PDF \cite{Monahan:2017hpu}. In this study, we present a systematic approach for calculating the matching from the gradient-flow scheme to the $\overline{\rm MS}$ scheme for Wilson-line operators, such as those defined in eq.~(\ref{eq:defqoperator}) and eq.~(\ref{eq:defgoperator}), specifically in the limit of small flow time. Our approach is based on the one-dimensional auxiliary-field formalism, and we provide comprehensive details of the one-loop order matching calculations.
For more recent applications of perturbative gradient flow, see refs.~\cite{Artz:2019bpr,Rizik:2020naq,Suzuki:2020zue,Harlander:2020duo,Boers:2020uqr,Shindler:2021bcx,Harlander:2022tgk,Mereghetti:2021nkt,Harlander:2022vgf,Brambilla:2021egm,Crosas:2023anw} and the references therein. 


We structure the remaining content of the paper as follows. In section \ref{sec:EFT}, we introduce the one-dimensional auxiliary-field formalism and elaborate on its application in the renormalization of Wilson-line operators. In section \ref{sec:introGF}, we introduce the gradient-flow formalism and implement the small flow-time expansion to the local current operators, yielding the core formulas for our matching calculations. Section \ref{sec:oneloopcalc} is dedicated to the computation of the matching coefficients up to one-loop order. Our summary and conclusions can be found in section \ref{sec:summary}. In appendix \ref{app:quarkquasi}, we present the detailed matching calculation with full flow time dependence for the hadronic matrix element that is related to quasi-PDFs at zero external momenta. We compare these findings with existing results from ref.~\cite{Monahan:2017hpu}, as well as with the results of section \ref{sec:oneloopcalc} in the small flow-time limit.


%%%%%%%%%%%%%%%%%%%%%%%%%%%%%%%%%%%%%%%%%%%
